\documentclass[11pt]{article}

\usepackage[margin=1.1in]{geometry}
\usepackage{amsmath, amssymb}
\usepackage{graphicx}
\usepackage{xcolor}
\usepackage[colorlinks=true, linkcolor=black, urlcolor=blue]{hyperref}
\usepackage{parskip}
\usepackage{enumitem}
\usepackage{mdframed}

\definecolor{boxbg}{HTML}{F4F6FA}
\definecolor{rule}{HTML}{C9D2E3}

\newmdenv[backgroundcolor=boxbg, linecolor=rule, linewidth=0.7pt,
          innertopmargin=8pt, innerbottommargin=8pt,
          innerleftmargin=11pt, innerrightmargin=11pt,
          skipabove=10pt, skipbelow=10pt]{keybox}

\title{\bfseries How GLIMPSE Works\\[4pt]
\large and the mathematics behind its map, ECLIPSE}
\author{Timothy Hsu \\ \small Aresty Summer Science Research Program}
\date{}

\begin{document}
\maketitle

\begin{keybox}
\textbf{Who this is for.} Anyone who wants to understand what GLIMPSE does and why,
without taking anything on faith. You need to be comfortable with the idea of an
average and an arrow; everything else is built up as we go. Mathematical readers can
skip the plain-language paragraphs -- the equations stand on their own.
\end{keybox}

\section{The problem}

A model decides something about you. You did not get the loan. The screening says you
are at risk of diabetes. Two questions follow immediately, and they are not the same
question:

\begin{enumerate}[leftmargin=*, itemsep=2pt]
  \item \textbf{Why me, and what would have to change?}
  \item \textbf{How sure is this, really?}
\end{enumerate}

Most explanation tools answer the first and quietly skip the second. GLIMPSE
(\emph{Geometric Lens for Inspecting the Multiplicity of Sets and Explanations})
answers both, on one picture, by drawing three things: where you stand, where the
decision flips, and \emph{how much that flip line moves} if you had fitted a slightly
different model.

That last point is the one people find surprising, so it is worth stating plainly
before any mathematics.

\begin{keybox}
\textbf{There is rarely one model.} Train a model on the same data with a slightly
different random seed, a slightly different regularisation, a slightly different
split, and you get a \emph{different} model that is \emph{just as accurate}. These
near-tied models frequently disagree about individual people -- one says approve,
another says decline, and no amount of test-set accuracy can tell you which is right.
This collection of equally-good models is called a \textbf{Rashomon set}, after the
film in which four witnesses describe the same event four irreconcilable ways.

If your explanation is drawn from only one of them, it is more confident than the
evidence supports. GLIMPSE draws the whole set.
\end{keybox}

\section{Part 1 --- the score}

Your information is a list of $d$ numbers. Glucose, BMI, age, and so on. Call that
list $x$. Every number is first put on a common scale, so that ``one unit'' means
``one standard deviation'' for every feature and none of them dominates just because
it happens to be measured in larger units.

The model attaches a weight to each feature and adds them up:
\[
  s(x) \;=\; w^\top x + b \;=\; w_1 x_1 + w_2 x_2 + \cdots + w_d x_d + b .
\]
This single number $s(x)$ is your \textbf{score}. A positive weight $w_j$ means that
feature pushes you toward the bad outcome; a negative one pushes away.

The score is turned into a probability by the \textbf{logistic} (or sigmoid) function,
which squashes any number onto the range $0$ to $1$:
\[
  p(y = 1 \mid x) \;=\; \sigma\big(s(x)\big), \qquad
  \sigma(z) = \frac{1}{1 + e^{-z}} .
\]
Score $0$ means probability exactly $1/2$ -- the fence. That is why the set of people
with $s(x) = 0$ is called the \textbf{decision boundary}.

Because the score is a plain weighted sum, the boundary is \emph{flat}: a line in two
dimensions, a plane in three, a \emph{hyperplane} in $d$. This flatness is what makes
everything later possible.

\begin{keybox}
\textbf{Why the sigmoid, and not some other S-shaped curve?}
It is not an arbitrary choice. Suppose each outcome group forms a Gaussian (bell-shaped)
cloud in feature space, and the two clouds have the same shape $\Sigma$, differing only
in where they sit ($\mu_0$ versus $\mu_1$). Write down Bayes' rule for
$p(y{=}1 \mid x)$. Each bell curve contributes a squared term
$-\tfrac12 (x - \mu_k)^\top \Sigma^{-1} (x - \mu_k)$; because the shape $\Sigma$ is
shared, the genuinely quadratic parts are \emph{identical} in both and cancel when you
subtract. What is left is exactly linear:
\[
  \log \frac{p(y{=}1 \mid x)}{p(y{=}0\mid x)} = w^\top x + b,
  \qquad w = \Sigma^{-1}(\mu_1 - \mu_0).
\]
So the logistic model is not an approximation under these assumptions -- it is the
\emph{exact} answer. This is what is meant by the remark that ``$p(y \mid x)$ comes
from a Gaussian.''

Two consequences matter for the rest of this document. First, the right direction to
point at is $\Sigma^{-1}(\mu_1 - \mu_0)$, \emph{not} the difference of the group
averages $\mu_1 - \mu_0$; you have to correct for how the features vary together
before you point. Second, $p(y\mid x)$ depends on $x$ \emph{only} through the single
number $w^\top x$. All of the information about the outcome lives along \textbf{one
direction}. Remember that; Section 6 depends on it.
\end{keybox}

\section{Part 2 --- the fan of equally-good models}

Fitting a model means choosing the parameters $\theta = (w, b)$ that make the
training data least surprising, measured by the average \emph{log-loss} $L(\theta)$.
The fitted $\hat\theta$ is the bottom of that valley.

Now ask a different question: not ``which parameters are best?'' but ``which parameters
are \emph{good enough}?'' Fix a small tolerance $\varepsilon$ and collect every model
whose loss is within $\varepsilon$ of the best:
\[
  R_\varepsilon \;=\; \{\, \theta \;:\; L(\theta) \le L(\hat\theta) + \varepsilon \,\}.
\]
This is the \textbf{Rashomon set}. GLIMPSE uses $\varepsilon = 0.03$.

$R_\varepsilon$ is an awkward shape, so we approximate it. Near the bottom of a smooth
valley, any loss function looks like a parabola -- this is just a second-order Taylor
expansion, and the first-order term vanishes \emph{because} we expanded at the minimum,
where the gradient is zero:
\[
  L(\theta) \;\approx\; L(\hat\theta)
  \;+\; \tfrac{1}{2}(\theta - \hat\theta)^\top H (\theta - \hat\theta),
\]
where $H$ is the matrix of second derivatives (the Hessian). Substituting gives a clean
shape:
\[
  R_\varepsilon \;\approx\;
  \big\{\, \theta : (\theta-\hat\theta)^\top Q\, (\theta-\hat\theta) \le 1 \,\big\},
  \qquad Q = \frac{H}{2\varepsilon}.
\]

That is the equation of an \textbf{ellipsoid} -- a stretched, tilted ball centred on
the fitted model. Directions in which the loss rises steeply are directions in which
the ellipsoid is thin: the data pinned those weights down. Directions in which the loss
is nearly flat are directions in which it is long: the data barely constrained them,
so models far apart there are all equally defensible.

\subsection*{How much do the good models disagree about \emph{you}?}

Every model in the ellipsoid gives you a score. Sweeping over all of them traces out an
\emph{interval} of scores rather than a single value, and the half-width of that
interval has a closed form. Writing $\tilde{x} = (x, 1)$ for your features with a $1$
appended to absorb the intercept:
\[
  \boxed{\;
  \mathrm{dis}(x) \;=\; \big\| \, Q^{-1/2}\, \tilde{x} \, \big\|
  \;}
\]
and the equally-good models place your score somewhere in
$\;\hat{s}(x) \pm \mathrm{dis}(x)$.

This one number is the heart of GLIMPSE. A small $\mathrm{dis}(x)$ means every
defensible model agrees about you and the decision is solid. A large $\mathrm{dis}(x)$
means your outcome is an artefact of which model happened to be trained -- and you
deserve to be told that.

It also gives \textbf{robust recourse}. The usual advice ``lower your BMI by $2$ and
you flip'' is advice that flips \emph{the one model that was fitted}. GLIMPSE instead
asks for a change that flips \emph{every} model in the ellipsoid: it requires
$\hat{s}(x) - \mathrm{dis}(x)$ to cross the boundary, not just $\hat s (x)$. That
advice costs more, and it actually works.

\section{Part 3 --- why a map is needed, and why it is hard}

Everything above lives in $d$ dimensions -- $8$ for the diabetes data, $57$ for the
spam data. Nobody can see that. To show it to a person, we need to flatten it onto a
sheet of paper, and \emph{any} flattening throws information away. The only real
question is \textbf{which} information.

We call the choice of flattening a \textbf{lens}. Formally a lens picks two directions
$a_1, a_2$ in feature space and reports where you fall along each:
\[
  z \;=\; \big( a_1^\top (x - \bar{x}),\;\; a_2^\top (x - \bar{x}) \big)
  \;\in\; \mathbb{R}^2 .
\]

Because this is just two weighted sums, it is \emph{affine}, and affine maps have three
properties GLIMPSE cannot do without:

\begin{enumerate}[leftmargin=*, itemsep=3pt]
  \item \textbf{Straight stays straight.} The flat decision boundary in $d$ dimensions
        appears as an honest straight \emph{line} on screen.
  \item \textbf{It runs backwards, exactly.} Click anywhere on the map and there is a
        formula returning the corresponding person -- their glucose, their BMI. Without
        this, ``what would I have to change?'' has no answer.
  \item \textbf{It handles strangers.} A new person who was not in the training data
        gets a position from the same two dot products. This is what makes the ``enter
        your own numbers'' box possible at all.
\end{enumerate}

Popular methods such as t-SNE and UMAP make prettier pictures but give up some or all
of these. t-SNE cannot place a person it has not already seen, which by itself
disqualifies it for an interactive tool.

\section{Part 4 --- what is wrong with PCA}

The default lens is Principal Component Analysis. PCA picks the two directions along
which the data is most \emph{spread out}. Concretely, it takes the covariance matrix
$\Sigma$ and uses its top two \textbf{eigenvectors}.

\begin{keybox}
\textbf{Eigenvectors, informally.} A symmetric matrix $M$ describes an amount of
``interest'' in every direction: point a unit arrow $v$ and the quantity $v^\top M v$
says how much of it lies that way. Almost every arrow gets tilted when you apply $M$,
but a few special ones do not -- they only get longer or shorter:
\[
  M v = \lambda v .
\]
These are the eigenvectors, and $\lambda$ is how much stretching each receives. The
eigenvector with the largest $\lambda$ is precisely the direction that maximises
$v^\top M v$, and the second-largest is the best direction at right angles to it. So
\emph{``take the top two eigenvectors''} always means the same thing: \emph{the flat
plane holding as much of the interest described by $M$ as any plane can.}
\end{keybox}

PCA is a perfectly good answer to the question it asks. The trouble is that it asks the
wrong question. PCA looks only at $\Sigma$ -- how the features vary. It is never told
that there is a model, a decision, or a Rashomon set. The direction the decision
actually depends on is $w$, and $w$ only lands in PCA's plane by luck.

The Gaussian argument in Part 1 sharpens this into something stronger than a
complaint. The direction that carries the outcome information is
$\Sigma^{-1}(\mu_1-\mu_0)$, which requires knowing $\mu_1 - \mu_0$ -- information PCA
never receives. PCA \emph{cannot} find the right axis except by accident.

\section{Part 5 --- ECLIPSE}

ECLIPSE starts from a simple observation: we do not want a plane that shows one thing
well. We want a plane that shows \emph{several} things tolerably. So write down each
thing worth seeing as its own matrix, then blend them.

\subsection*{The four things worth seeing}

Each is a $d \times d$ symmetric matrix, and in each case the quantity $v^\top S_i v$
measures ``how much of this you would see along direction $v$.''

\begin{enumerate}[leftmargin=*, itemsep=5pt]
  \item \textbf{Spread} --- $S_0 = \mathrm{Cov}(X)$. \\
        People should not all land on top of each other. This is exactly PCA's
        criterion, kept as one voice among four rather than the only one.

  \item \textbf{The model} --- $S_1 = w w^\top$. \\
        A rank-one matrix built from the weight vector. Since
        $v^\top (w w^\top) v = (v^\top w)^2$, it is large exactly when $v$ points along
        $w$. Including it pulls the map toward showing the score as a left-to-right
        position.

  \item \textbf{Disagreement} --- $S_2 = \mathbb{E}\!\left[\nabla \mathrm{dis}\,
        \nabla \mathrm{dis}^\top\right]$. \\
        $\nabla \mathrm{dis}$ is the direction in which the disagreement
        $\mathrm{dis}(x)$ changes fastest. Averaging the outer product over everybody
        gives the directions along which model uncertainty varies most. Keeping these
        in the plane is what lets a user \emph{see} where the good models stop agreeing.

  \item \textbf{Imputation} --- $S_3 = \mathrm{diag}(f_j \cdot \sigma_j^2)$. \\
        $f_j$ is the fraction of people whose feature $j$ was missing and had to be
        filled in. This keeps guessed-at features visible instead of hiding the fact
        that part of the picture was invented.
\end{enumerate}

Each matrix is rescaled to have trace $1$, so that ``one unit of interest'' means the
same thing in all four and the weights below are comparable.

\subsection*{The blend}

Choose four non-negative weights summing to one and add the matrices up:
\[
  \boxed{\;
    M \;=\; \omega_0 S_0 \;+\; \omega_1 S_1 \;+\; \omega_2 S_2 \;+\; \omega_3 S_3,
    \qquad \textstyle\sum_i \omega_i = 1, \;\; \omega_i \ge 0.
  \;}
\]
Then take the top two eigenvectors of $M$ as the two axes of the map.

That is the entire method. Its justification is the fact in the eigenvector box: the
top-two eigenvectors of $M$ give the plane maximising total interest $v^\top M v$, so
ECLIPSE returns \emph{the flat plane that best satisfies all four demands at once, in
the proportions you asked for}. Formally it maximises the sum of the two largest values
of $A^\top M A$ over planes $A$, which is a Ky Fan extremal problem with exactly this
solution.

Two sanity checks:
\begin{itemize}[leftmargin=*, itemsep=2pt]
  \item $\omega = (1, 0, 0, 0)$ gives $M = \mathrm{Cov}(X)$ --- \textbf{ECLIPSE reduces
        exactly to PCA}. PCA is a special case, not a rival.
  \item $\omega = (0, 1, 0, 0)$ gives $M = ww^\top$, which has rank one: the first axis
        becomes $w$ itself and the second is unconstrained. Showing the score perfectly
        and nothing else.
\end{itemize}
Every useful lens sits between these two extremes, which is what the weights are for.

\subsection*{Choosing the weights without cheating}

Hand-tuning the weights per dataset would be a way of fitting the picture to whatever
we hoped to see. ECLIPSE instead chooses them by a rule, on data held out for the
purpose:

\begin{enumerate}[leftmargin=*, itemsep=3pt]
  \item Split off a validation set the lens is not fitted on.
  \item Compute PCA's neighbourhood faithfulness (\emph{trustworthiness}) there. This
        is the benchmark.
  \item Sweep candidate weightings. \textbf{Reject} any whose trustworthiness falls
        more than $\tau = 0.05$ below PCA's. This is the promise that the map never
        becomes much more distorted than the familiar one.
  \item Among the survivors, keep the one that best preserves the score ordering
        \emph{and} the disagreement ordering (the product of the two rank
        correlations, so a candidate must do acceptably at both).
\end{enumerate}

The result: at most $0.05$ of neighbourhood faithfulness is ever spent, and it is spent
on the two things a decision map exists to show.

\section{Part 6 --- what the Gaussian argument buys}

Return to the fact from Part 1: under Gaussian outcome groups with a shared shape,
$p(y \mid x)$ depends on $x$ only through $w^\top x$. The score is a \textbf{sufficient
statistic}.

This has a sharp design consequence. If \emph{all} the outcome information lies along
one direction, then a second ``outcome axis'' on the map is not merely unhelpful --
there is provably nothing left for it to carry. The correct architecture for a 2-D
decision map is therefore:

\begin{center}
\emph{axis 1 carries the score; axis 2 carries something the score cannot tell you.}
\end{center}

Which is precisely the shape of the blend: $S_1$ claims the first axis, and $S_0$,
$S_2$, $S_3$ compete for the second. It also explains something that looks like a flaw
elsewhere: Linear Discriminant Analysis yields only \emph{one} discriminant direction
for a two-class problem. That is not a limitation of LDA. It is LDA being right about a
problem that has only one informative direction.

Taking the argument literally suggests a variant -- set $a_1 = w / \|w\|$ exactly, and
give $a_2$ the best of what remains. Tested across seven datasets, this makes the score
a \emph{perfect} function of the horizontal position (rank correlation $1.000$ on all
seven, against $0.997$ for the tuned rule), confirming the theory numerically.

\begin{keybox}
\textbf{Two different Gaussians.} They are easy to confuse.
\begin{itemize}[leftmargin=*, itemsep=1pt]
  \item A Gaussian over \textbf{$x$ given $y$} (Part 1) makes the log-odds linear, so
        the boundary is a straight line.
  \item A Gaussian over \textbf{$\theta$ given the data} (Part 2) makes the Rashomon
        set an ellipsoid --- the parabola approximation there is the same thing as a
        Gaussian approximation to the posterior over parameters.
\end{itemize}
One assumption makes the boundary a line; the other makes the uncertainty an ellipse.
GLIMPSE draws exactly those two objects.
\end{keybox}

\section{Part 7 --- knowing in advance when this cannot work}

An honest method should be able to say when it will fail. ECLIPSE can.

No flat plane can show everything. For the disagreement term specifically there is a
hard ceiling. Let $\alpha_1 \ge \alpha_2 \ge \cdots \ge \alpha_d$ be the eigenvalues of
$\Sigma^{1/2} A\, \Sigma^{1/2}$, where $A$ is the matrix defining $\mathrm{dis}$. Then
define
\[
  \kappa_2 \;=\; \sqrt{ \frac{\alpha_1^2 + \alpha_2^2}
                             {\alpha_1^2 + \cdots + \alpha_d^2} } \;\in\; [0, 1].
\]
$\kappa_2$ bounds how well \emph{any} two-dimensional plane can represent model
disagreement -- not just ECLIPSE's, any plane at all.

The useful part is that $\kappa_2$ is computable \emph{before fitting any lens}. If it
is close to $1$, the disagreement structure is essentially two-dimensional and a map
can show it. If it is small, that structure is spread across many directions and no
picture will capture it; weight spent on $S_2$ is wasted and should go to the score
instead.

Measured across seven datasets this behaves as advertised: $\kappa_2$ ranges from
$0.37$ (German credit) to $0.96$ (breast cancer), and it correctly separates the
datasets where the disagreement axis pays for itself from those where it cannot.
Using $\kappa_2$ directly in place of the validation sweep gives statistically
indistinguishable maps roughly $170\times$ faster.

\section{Summary}

\begin{keybox}
\begin{enumerate}[leftmargin=*, itemsep=4pt]
  \item Your information is a list of numbers $x$; the model scores you with a weighted
        sum $s(x) = w^\top x + b$ and the sigmoid turns that into a probability. Under
        Gaussian outcome groups, this is not an approximation but the exact posterior.
  \item Many \emph{different} models fit the data equally well. To second order they
        form an ellipsoid around the fitted one, and the spread of scores they assign
        you is $\mathrm{dis}(x) = \|Q^{-1/2}\tilde{x}\|$ --- how much your decision
        depends on which model happened to be trained.
  \item Showing this to a person means flattening $d$ dimensions onto two. GLIMPSE uses
        an affine lens, which keeps straight boundaries straight, runs backwards
        exactly, and can place a person it has never seen.
  \item PCA chooses that plane using spread alone and is never told a decision exists.
  \item ECLIPSE writes each thing worth seeing as a matrix --- spread, model,
        disagreement, imputation --- blends them with weights $\omega$, and takes the
        top two eigenvectors of the blend. $\omega = (1,0,0,0)$ recovers PCA exactly.
  \item The weights are chosen by a rule on held-out data, never by hand, and the rule
        may never give up more than $\tau = 0.05$ of neighbourhood faithfulness
        relative to PCA.
  \item $\kappa_2$ says in advance how much of the disagreement structure \emph{any}
        flat map could possibly show, so the method knows when to stop trying.
\end{enumerate}
\end{keybox}

\vspace{6pt}
\noindent\textbf{Notation.}\quad
$x$ features (standardised) $\cdot$
$y$ outcome $\cdot$
$w, b$ weights and intercept $\cdot$
$\theta = (w,b)$ $\cdot$
$\sigma(\cdot)$ logistic function $\cdot$
$\Sigma$ covariance of the features $\cdot$
$H$ Hessian of the loss $\cdot$
$\varepsilon$ Rashomon tolerance $\cdot$
$Q = H/2\varepsilon$ $\cdot$
$\mathrm{dis}(x)$ disagreement half-width $\cdot$
$S_0 \ldots S_3$ interest matrices $\cdot$
$\omega$ blend weights $\cdot$
$\tau$ trustworthiness budget $\cdot$
$\kappa_2$ disagreement ceiling.

\end{document}
